\begin{helper}[Reconstruction at a maximal restricted endpoint]
Let corresponding ETTs on $J$ and $U$ be given for the same original ETT
and coloring $\sigma$.  Suppose $J$ carries a hierarchy base $B_J$, a
hierarchy $K_J$, and a good witness $\mathcal G_J$, and suppose there is a
$\sigma$-\noderef{TashkinovAugmentationSequence} from $J$ to $U$.  Put
$d=K_J.\mathit{dividerCount}$, $A=K_J.\mathit{starred}(d)$,
$D=D^{K_J}_{n,d}$, and $\Gamma_\eta=\mathcal G_J.\Gamma^d_\eta$.  Assume
that $\Gamma_\eta$ avoids the last-stage segment colors of $U$ for every
$\eta\in D$, and that no one-step augmentation of $U$ satisfies all these
same fixed-Gamma avoidance conditions.

Then the corresponding ETT on $U$ has a hierarchy base $B_U$ with the
same core, closure, and numerical supporting position as $B_J$, and a
hierarchy $K_U$ with $d$ dividers.  Under the canonical index
identification, all its nonterminal stages, starred stages, indexed stored
sets, stage-missing sets, divider positions, and Gamma pairs agree with
those of $K_J$.  Its terminal stored set is contained in $D$, and $U$ is
minus-closed outside the union of the terminal Gamma pairs.
\end{helper}
\begin{proof}
We first construct the correspondence needed to transport the hierarchy
base.  Write $\mathcal E_J$ and $\mathcal E_U$ for the replacement ETTs
in the certificates on $J$ and $U$.  Their rung-agreement fields identify
both rung numbers with that of the original ETT.  After transporting every
dependent index across those equalities, their ladder-tree, raw stored-set,
connecting-edge, and extension-kind agreements in
\noderef{CorrespondingExtendedTashkinovTree} compose through the common
original series.  The series-presentation and final-color fields likewise
give the required presentation of $\mathcal E_U$ and the same displayed
final coloring $\sigma$ when $\mathcal E_J$ is used as the base.

The stability field needs a separate rung-by-rung argument.  Fix a
positive-stage index $i$, and denote the three colorings there by
$\varphi_i$ in the original series, $\psi_i$ in the series on $J$, and
$\chi_i$ in the series on $U$.  The two correspondence certificates give
that $\psi_i$ and $\chi_i$ are stable relative to $\varphi_i$ on the
common exact ladder tree and original protected set.  They also state that
$\psi_i$ and $\chi_i$ are proper and leave exactly $e_0$ uncolored.  The
certified-step field of \noderef{TashkinovSeries} supplies a positive
\noderef{TashkinovSeriesStep} constructor for $\varphi_i$.  In the SE
constructor, its coloring is the selected coloring of
\noderef{MaximumDefectiveChoice}, whose fields give properness and the
one-uncolored property; the PE constructor records those two properties
directly; and the RE constructor identifies the current coloring with the
proper, one-uncolored preceding coloring.  Thus all three colorings define
elements of the subtype occurring in \noderef{sc2}.

Apply symmetry in \noderef{sc2} to reverse the stability from
$\varphi_i$ to $\psi_i$, and then apply transitivity there with the
stability from $\varphi_i$ to $\chi_i$.  This gives stability from
$\psi_i$ to $\chi_i$.  To put that relation in precisely the indices and
protected set required over $\mathcal E_J$, use the common-ladder-tree
equality above and unfold \noderef{TashkinovStoredColors}.  The raw stored
entries agree with the original ones.  Moreover, the vertexwise
missing-set clause of \noderef{StableColoring} for $\psi_i$ identifies the
tree-missing sets of $\psi_i$ and $\varphi_i$ on that common ladder tree.
Consequently the computed stored set for the series on $J$ is the original
protected set.  Transporting the stability relation across these two
equalities gives exactly the replacement-stability field from
$\mathcal E_J$ to $\mathcal E_U$.  Doing this for every positive-stage
index, and using the proper and one-uncolored fields from the certificate
on $U$, completes an auxiliary
\noderef{CorrespondingExtendedTashkinovTree} certificate with base
$\mathcal E_J$ and replacement ETT $\mathcal E_U$.

Apply \noderef{HierarchyBaseCorrespondingTransport} to this auxiliary
certificate and $B_J$.  It gives a base $B_U$ on $\mathcal E_U$ with
exactly the same core, closure, and numerical supporting position.  This
also handles the PE support vertex: the cited transport lemma compares the
exact chosen boundary-edge identities and their inside endpoints, rather
than assuming uniqueness of a support vertex.

We next fix the prefix fact used throughout the reconstruction.  By
\noderef{TashkinovAugmentationSequenceInitialSegment}, the given
$\sigma$-\noderef{TashkinovAugmentationSequence} supplies
\[
                         J\preccurlyeq U.                       \tag{1}
\]
Thus the length clause of \noderef{TreeSequenceInitialSegment} gives
$|J|\leq|U|$, and its vertex- and edge-entry clauses identify all entries
through $|J|$.

Construct $K_U$ with divider count $d$, copy stages $0,\ldots,d$ from
$K_J$, and take $U$ as stage $d+1$.  For nesting and strictness before the
last step, the corresponding fields of
\noderef{TashkinovHierarchy} copy verbatim.  At the last step, its old
nesting gives the copied stage $d$ as an exact prefix of the old terminal
stage $J$; composing its length and entrywise clauses with (1) gives the
required prefix of $U$.  The old strict inequality into $J$, followed by
$|J|\leq|U|$, gives strictness into $U$.  Embed every old divider position
into $\operatorname{Fin}(|U|+1)$ using (1).  The numerical values,
positivity, and strict increase do not change.

For an old divider position $p$, the divider-stage field of
\noderef{TashkinovHierarchy} identifies its stage with
$\noderef{TreeSequencePrefix}(J,p)$.  Since $p\leq|J|$ and (1) identifies
the entries of $J$ and $U$ through $p$, the inherited vertex and edge maps
in the two prefix structures agree entry by entry; proof irrelevance
identifies their attachment and injectivity witnesses.  Hence this same
stage is $\noderef{TreeSequencePrefix}(U,p)$, as required.  If $d>0$, the
core-inside-first field has the same copied target stage and therefore
copies.  If $d=0$, that target is the new terminal stage $U$; compose the
old core-to-$J$ prefix with (1), using the reconstructed core equality.
This verifies every field of $K_U$.  It also proves the asserted stage and
numerical divider equalities.  Substitution in the definitions of
\noderef{HierarchyStarredStage},
\noderef{HierarchyStoredColorsAt}, and
\noderef{HierarchyStageMissingColors} then give the remaining indexed
equalities in the conclusion.

Define the Gamma pairs of $\mathcal G_U$ by copying those of
$\mathcal G_J$ under this index identification.  The pair-cardinality,
pairwise-disjointness, unique-refresh, PE-base, and internal minus-closed
fields of \noderef{GoodHierarchy} use only the copied base, stages, stored
sets, stage-missing sets, starred stages, and Gamma pairs, so each is the
corresponding field of $\mathcal G_J$ after the displayed equalities.

It remains to verify every avoidance field, because
\noderef{HierarchySegmentColors} also depends on the whole terminal tree.
Fix $j<d$, let $B$ be the common copied next stage, and fix a color
$\eta$.  Iterating the nestedness fields of
\noderef{TashkinovHierarchy} gives $B\preccurlyeq J$, so $|B|\leq|J|$.
Apply \noderef{HierarchyFirstMissingPositionPrefixCutoff} to (1) with
cutoff $u=|B|$.  It gives
\[
 \min\{p_J(\eta),|B|\}=\min\{p_U(\eta),|B|\}.              \tag{2}
\]
The lower starred stage and the edge map of $B$ are unchanged.  Therefore,
for every edge index of $B$, the lower-bound test, the strict upper-bound
test in (2), and the edge-color test in the defining filter of
\noderef{HierarchySegmentColors} agree entry by entry.  The two segment
color sets are equal.  Notice that the cutoff theorem proves (2) also when
$\eta$ is absent from $J$ or $U$, using the totalized last-position
fallback; no missing-color witness is being assumed.  Hence the old
avoidance field transfers for every $j<d$.

At $j=d$, the starred stage is the fixed tree $A$, the next stage and
terminal stage are both $U$, and the segment filter is exactly the
restricted-color set appearing in the hypothesis.  The old avoidance
field still supplies containment of $\Gamma_\eta$ in the unchanged
stage-missing set, while the endpoint hypothesis supplies disjointness
from this new segment set.  Together they are precisely the last avoidance
field of \noderef{GoodHierarchy}.  Thus $\mathcal G_U$ is a good witness,
and its definition gives every claimed Gamma equality.

We now compare terminal stored sets.  Let $R$ be the computed terminal-rung
\noderef{TashkinovStoredColors} set for the series on $J$.  The
rung-by-rung stored-set calculation above identifies it with the
corresponding set for the series on $U$.  The last starred stage $A$ is an
exact prefix of $J$: when $d=0$ this is the core-inside-first field, and
when $d>0$ it is the last nestedness field, after unfolding
\noderef{HierarchyStarredStage}.  Compose that prefix with (1).  The image
description in \noderef{TreeSequenceVertexSet} then gives
$V(A)\subseteq V(U)$, and the union definition in
\noderef{TreeMissingColors} gives
$\overline\sigma(A)\subseteq\overline\sigma(U)$.  Unfolding
\noderef{HierarchyStoredColorsAt} and
\noderef{HierarchyTerminalStoredColors} now gives
\[
 D^{K_U}_{n,d+1}=R\setminus\overline\sigma(U),\qquad
 D^{K_J}_{n,d}=R\setminus\overline\sigma(A).                \tag{3}
\]
The missing-set inclusion proves, element by element, that the first set
is contained in the second, as claimed.

Finally suppose the desired \noderef{TreeColorSetMinusClosed} conclusion
fails.  By \noderef{TreeColorSetClosed}, there are a color $c$ missing on
$U$ but outside the union of Gamma pairs indexed by
$D^{K_U}_{n,d+1}$, and an exact identity
$f\in\noderef{BoundaryColorEdges}(G,\sigma,V(U),c)$.  By
\noderef{BoundaryMissingColorAugmentation}, adjoining $f$ gives a tree
$W$ with an exact relation
\noderef{TreeSequenceAdjoinEdge}$(U,f,W)$ and a one-step
$\sigma$-\noderef{TashkinovAugmentation} from $U$ to $W$.  In particular,
\[
 U\preccurlyeq W,\qquad |W|=|U|+1,
 \qquad W.\mathit{edge}(|U|)=f,\qquad \sigma(f)=c.           \tag{4}
\]

Fix $\eta\in D=D^{K_J}_{n,d}$.  First suppose
$\eta\in D^{K_U}_{n,d+1}$.  The choice of $c$ gives
$c\notin\Gamma_\eta$.  Apply
\noderef{HierarchyFirstMissingPositionPrefixCutoff} to the prefix in (4)
with cutoff $|U|$.  If an edge index contributing to the restricted segment
of $W$ is below $|U|$, the cutoff equality and the entrywise part of (4)
show that its color already belongs to the restricted segment of $U$.
The only remaining possible index is $|U|$, by the one-more length in (4),
and its color is $c$.  Thus the new segment-color set is contained in the
old one together with $\{c\}$.  Endpoint avoidance and
$c\notin\Gamma_\eta$ prove the required disjointness.  This argument also
covers the no-missing fallback, because the cited cutoff theorem was
applied before separating the possible new last index.

Now suppose $\eta\notin D^{K_U}_{n,d+1}$.  Formula (3) and
$\eta\in D$ imply that $\eta$ is missing at an actual vertex of $U$.
By (4), every old vertex entry is unchanged in $W$.  The least such
missing vertex, and the absence of a missing witness at every earlier
entry, therefore show directly from
\noderef{HierarchyFirstMissingPosition} that the first missing position is
the same in $U$ and $W$, even when that witness is the last vertex of $U$.
The new edge has index $|U|$, which is not strictly before that position;
all earlier edge entries are unchanged.  Hence the restricted segment set
for $W$ equals that for $U$, so endpoint avoidance again persists.

Thus $W$ satisfies every fixed-Gamma avoidance condition.  Its one-step
augmentation contradicts the assumed endpoint maximality.  No such $c$
and $f$ exist, and unfolding \noderef{TreeColorSetMinusClosed} now gives
the asserted terminal minus-closedness.
\end{proof}
